% figs_part3.tex —— 由 mit/make_figs_part3.py 產生，請勿手動修改。
% 重新產生：python mit/make_figs_part3.py

\newcommand{\FigMnineStuck}{%
\begin{tikzpicture}
\fill[cA!14] (-0.25,1.25) -- (0.25,1.25) -- (0.25,1.75) -- (-0.25,1.75) -- cycle;
\draw[cA,line width=0.8pt] (-0.25,1.25) -- (0.25,1.25) -- (0.25,1.75) -- (-0.25,1.75) -- cycle;
\node[font=\small] at (0,1.5) {$1$};
\draw[gray!55,dashed,line width=0.5pt] (0.35,1.25) -- (0.85,1.25) -- (0.85,1.75) -- (0.35,1.75) -- cycle;
\draw[gray!55,dashed,line width=0.5pt] (0.95,1.25) -- (1.45,1.25) -- (1.45,1.75) -- (0.95,1.75) -- cycle;
\fill[cB!18] (2.3,-0.4) -- (2.8,-0.4) -- (2.8,0.1) -- (2.3,0.1) -- cycle;
\draw[cB,line width=1.1pt] (2.3,-0.4) -- (2.8,-0.4) -- (2.8,0.1) -- (2.3,0.1) -- cycle;
\node[font=\small] at (2.55,-0.15) {$2$};
\fill[cBad!22] (2.3,0.2) -- (2.8,0.2) -- (2.8,0.7) -- (2.3,0.7) -- cycle;
\draw[cBad,line width=1.1pt] (2.3,0.2) -- (2.8,0.2) -- (2.8,0.7) -- (2.3,0.7) -- cycle;
\node[font=\small] at (2.55,0.45) {$3$};
\draw[gray!55,dashed,line width=0.5pt] (3.7,1.25) -- (4.2,1.25) -- (4.2,1.75) -- (3.7,1.75) -- cycle;
\draw[gray!55,dashed,line width=0.5pt] (4.3,1.25) -- (4.8,1.25) -- (4.8,1.75) -- (4.3,1.75) -- cycle;
\draw[gray!55,dashed,line width=0.5pt] (4.9,1.25) -- (5.4,1.25) -- (5.4,1.75) -- (4.9,1.75) -- cycle;
\node[font=\small,text=gray!40!black] at (0.6,2.02) {輸出};
\node[font=\small,text=gray!40!black] at (4.55,2.02) {輸入（已用完）};
\node[font=\small,text=gray!40!black] at (2.55,-0.75) {堆疊};
\draw[gray!60!black,line width=1.2pt] (2.19,1.05) -- (2.19,-0.45) -- (2.91,-0.45) -- (2.91,1.05);
\draw[-{Stealth[length=2.2mm]},gray!60!black,line width=0.7pt] (3.6,1.5) -- (2.7,1.17);
\draw[-{Stealth[length=2.2mm]},gray!60!black,line width=0.7pt] (2.4,1.17) -- (1.55,1.5);
\node[font=\footnotesize,text=gray!40!black] at (3.35,1) {推入};
\node[font=\footnotesize,text=gray!40!black] at (1.75,1) {彈出};
\node[font=\small] at (2.55,-1.25) {(a) 堆疊排 $231$};
\fill[cA!14] (6.85,0.1) -- (7.35,0.1) -- (7.35,0.6) -- (6.85,0.6) -- cycle;
\draw[cA,line width=0.8pt] (6.85,0.1) -- (7.35,0.1) -- (7.35,0.6) -- (6.85,0.6) -- cycle;
\node[font=\small] at (7.1,0.35) {$1$};
\draw[gray!55,dashed,line width=0.5pt] (7.45,0.1) -- (7.95,0.1) -- (7.95,0.6) -- (7.45,0.6) -- cycle;
\draw[gray!55,dashed,line width=0.5pt] (8.05,0.1) -- (8.55,0.1) -- (8.55,0.6) -- (8.05,0.6) -- cycle;
\fill[cBad!22] (9.22,0.1) -- (9.72,0.1) -- (9.72,0.6) -- (9.22,0.6) -- cycle;
\draw[cBad,line width=1.1pt] (9.22,0.1) -- (9.72,0.1) -- (9.72,0.6) -- (9.22,0.6) -- cycle;
\node[font=\small] at (9.47,0.35) {$3$};
\fill[cB!18] (9.82,0.1) -- (10.32,0.1) -- (10.32,0.6) -- (9.82,0.6) -- cycle;
\draw[cB,line width=1.1pt] (9.82,0.1) -- (10.32,0.1) -- (10.32,0.6) -- (9.82,0.6) -- cycle;
\node[font=\small] at (10.07,0.35) {$2$};
\draw[gray!55,dashed,line width=0.5pt] (11.1,0.1) -- (11.6,0.1) -- (11.6,0.6) -- (11.1,0.6) -- cycle;
\draw[gray!55,dashed,line width=0.5pt] (11.7,0.1) -- (12.2,0.1) -- (12.2,0.6) -- (11.7,0.6) -- cycle;
\draw[gray!55,dashed,line width=0.5pt] (12.3,0.1) -- (12.8,0.1) -- (12.8,0.6) -- (12.3,0.6) -- cycle;
\node[font=\small,text=gray!40!black] at (7.7,-0.27) {輸出};
\node[font=\small,text=gray!40!black] at (11.95,-0.27) {輸入（已用完）};
\node[font=\small,text=gray!40!black] at (9.825,-0.27) {佇列};
\node[font=\footnotesize,text=gray!40!black] at (9.3,0.9) {頭};
\node[font=\footnotesize,text=gray!40!black] at (10.35,0.9) {尾};
\draw[gray!60!black,line width=1.2pt] (9.05,0.71) -- (10.6,0.71);
\draw[gray!60!black,line width=1.2pt] (9.05,-0.01) -- (10.6,-0.01);
\draw[-{Stealth[length=2.2mm]},gray!60!black,line width=0.7pt] (11.05,0.35) -- (10.65,0.35);
\draw[-{Stealth[length=2.2mm]},gray!60!black,line width=0.7pt] (9,0.35) -- (8.65,0.35);
\draw[cB,dashed,line width=0.9pt,-{Stealth[length=2.2mm]}] (11.35,0.68) -- (11.35,1.5) -- (8.3,1.5) -- (8.3,0.68);
\node[font=\small,text=cB!80!black] at (9.825,1.75) {插隊線};
\node[font=\small] at (9.85,-1.25) {(b) 佇列＋插隊線排 $321$};
\end{tikzpicture}}

\newcommand{\FigMtenPlots}{%
\begin{tikzpicture}
\draw[gray!45,line width=0.4pt] (0.162,0.162) -- (1.278,0.162) -- (1.278,1.278) -- (0.162,1.278) -- cycle;
\draw[gray!18,line width=0.3pt] (0.36,0.162) -- (0.36,1.278);
\draw[gray!18,line width=0.3pt] (0.162,0.36) -- (1.278,0.36);
\draw[gray!18,line width=0.3pt] (0.72,0.162) -- (0.72,1.278);
\draw[gray!18,line width=0.3pt] (0.162,0.72) -- (1.278,0.72);
\draw[gray!18,line width=0.3pt] (1.08,0.162) -- (1.08,1.278);
\draw[gray!18,line width=0.3pt] (0.162,1.08) -- (1.278,1.08);
\draw[cTop,line width=0.9pt] (0.36,0.72) -- (0.72,1.08);
\draw[cTop,line width=0.9pt] (0.72,1.08) -- (1.08,0.36);
\fill[cTop] (0.36,0.72) circle (0.094);
\node[font=\small,text=cTop] at (0.36,-0.058) {$2$};
\fill[cTop] (0.72,1.08) circle (0.094);
\node[font=\small,text=cTop] at (0.72,-0.058) {$3$};
\fill[cTop] (1.08,0.36) circle (0.094);
\node[font=\small,text=cTop] at (1.08,-0.058) {$1$};
\node[font=\small] at (0.72,-0.55) {模式 $\pi=231$};
\draw[gray!45,line width=0.4pt] (3.112,0.162) -- (4.948,0.162) -- (4.948,1.998) -- (3.112,1.998) -- cycle;
\draw[gray!18,line width=0.3pt] (3.31,0.162) -- (3.31,1.998);
\draw[gray!18,line width=0.3pt] (3.112,0.36) -- (4.948,0.36);
\draw[gray!18,line width=0.3pt] (3.67,0.162) -- (3.67,1.998);
\draw[gray!18,line width=0.3pt] (3.112,0.72) -- (4.948,0.72);
\draw[gray!18,line width=0.3pt] (4.03,0.162) -- (4.03,1.998);
\draw[gray!18,line width=0.3pt] (3.112,1.08) -- (4.948,1.08);
\draw[gray!18,line width=0.3pt] (4.39,0.162) -- (4.39,1.998);
\draw[gray!18,line width=0.3pt] (3.112,1.44) -- (4.948,1.44);
\draw[gray!18,line width=0.3pt] (4.75,0.162) -- (4.75,1.998);
\draw[gray!18,line width=0.3pt] (3.112,1.8) -- (4.948,1.8);
\draw[cTop,line width=0.9pt] (3.31,1.44) -- (4.03,1.8);
\draw[cTop,line width=0.9pt] (4.03,1.8) -- (4.39,1.08);
\fill[cTop] (3.31,1.44) circle (0.094);
\node[font=\small,text=cTop] at (3.31,-0.058) {$4$};
\fill[gray!55] (3.67,0.36) circle (0.075);
\node[font=\small,text=black] at (3.67,-0.058) {$1$};
\fill[cTop] (4.03,1.8) circle (0.094);
\node[font=\small,text=cTop] at (4.03,-0.058) {$5$};
\fill[cTop] (4.39,1.08) circle (0.094);
\node[font=\small,text=cTop] at (4.39,-0.058) {$3$};
\fill[gray!55] (4.75,0.72) circle (0.075);
\node[font=\small,text=black] at (4.75,-0.058) {$2$};
\node[font=\small] at (4.03,-0.55) {$w=41532$：含有 $231$};
\draw[gray!45,line width=0.4pt] (6.962,0.162) -- (8.798,0.162) -- (8.798,1.998) -- (6.962,1.998) -- cycle;
\draw[gray!18,line width=0.3pt] (7.16,0.162) -- (7.16,1.998);
\draw[gray!18,line width=0.3pt] (6.962,0.36) -- (8.798,0.36);
\draw[gray!18,line width=0.3pt] (7.52,0.162) -- (7.52,1.998);
\draw[gray!18,line width=0.3pt] (6.962,0.72) -- (8.798,0.72);
\draw[gray!18,line width=0.3pt] (7.88,0.162) -- (7.88,1.998);
\draw[gray!18,line width=0.3pt] (6.962,1.08) -- (8.798,1.08);
\draw[gray!18,line width=0.3pt] (8.24,0.162) -- (8.24,1.998);
\draw[gray!18,line width=0.3pt] (6.962,1.44) -- (8.798,1.44);
\draw[gray!18,line width=0.3pt] (8.6,0.162) -- (8.6,1.998);
\draw[gray!18,line width=0.3pt] (6.962,1.8) -- (8.798,1.8);
\fill[cA] (7.16,1.8) circle (0.075);
\node[font=\small,text=black] at (7.16,-0.058) {$5$};
\fill[cA] (7.52,0.72) circle (0.075);
\node[font=\small,text=black] at (7.52,-0.058) {$2$};
\fill[cA] (7.88,0.36) circle (0.075);
\node[font=\small,text=black] at (7.88,-0.058) {$1$};
\fill[cA] (8.24,1.44) circle (0.075);
\node[font=\small,text=black] at (8.24,-0.058) {$4$};
\fill[cA] (8.6,1.08) circle (0.075);
\node[font=\small,text=black] at (8.6,-0.058) {$3$};
\node[font=\small] at (7.88,-0.55) {$w=52143$：避開 $231$};
\end{tikzpicture}}

\newcommand{\FigMtenGrid}{%
\begin{tikzpicture}
\fill[cA!12,rounded corners=2pt] (1.22,-0.19) rectangle (2.18,0.19);
\node[font=\small] at (1.7,0) {${\color{cA}\mathbf{123}}$};
\fill[cA!12,rounded corners=2pt] (2.52,-0.19) rectangle (3.48,0.19);
\node[font=\small] at (3,0) {${\color{cA}\mathbf{132}}$};
\fill[cA!12,rounded corners=2pt] (3.82,-0.19) rectangle (4.78,0.19);
\node[font=\small] at (4.3,0) {${\color{cA}\mathbf{213}}$};
\node[font=\small] at (5.6,0) {${\color{cBad}2}{\color{cBad}3}{\color{cBad}1}$};
\fill[cA!12,rounded corners=2pt] (6.42,-0.19) rectangle (7.38,0.19);
\node[font=\small] at (6.9,0) {${\color{cA}\mathbf{312}}$};
\fill[cA!12,rounded corners=2pt] (7.72,-0.19) rectangle (8.68,0.19);
\node[font=\small] at (8.2,0) {${\color{cA}\mathbf{321}}$};
\fill[cA!12,rounded corners=2pt] (1.22,-0.87) rectangle (2.18,-0.49);
\node[font=\small] at (1.7,-0.68) {${\color{cA}\mathbf{1234}}$};
\fill[cA!12,rounded corners=2pt] (2.52,-0.87) rectangle (3.48,-0.49);
\node[font=\small] at (3,-0.68) {${\color{cA}\mathbf{1243}}$};
\fill[cA!12,rounded corners=2pt] (3.82,-0.87) rectangle (4.78,-0.49);
\node[font=\small] at (4.3,-0.68) {${\color{cA}\mathbf{1324}}$};
\node[font=\small] at (5.6,-0.68) {${\color{gray}1}{\color{cBad}3}{\color{cBad}4}{\color{cBad}2}$};
\fill[cA!12,rounded corners=2pt] (6.42,-0.87) rectangle (7.38,-0.49);
\node[font=\small] at (6.9,-0.68) {${\color{cA}\mathbf{1423}}$};
\fill[cA!12,rounded corners=2pt] (7.72,-0.87) rectangle (8.68,-0.49);
\node[font=\small] at (8.2,-0.68) {${\color{cA}\mathbf{1432}}$};
\fill[cA!12,rounded corners=2pt] (1.22,-1.37) rectangle (2.18,-0.99);
\node[font=\small] at (1.7,-1.18) {${\color{cA}\mathbf{2134}}$};
\fill[cA!12,rounded corners=2pt] (2.52,-1.37) rectangle (3.48,-0.99);
\node[font=\small] at (3,-1.18) {${\color{cA}\mathbf{2143}}$};
\node[font=\small] at (4.3,-1.18) {${\color{cBad}2}{\color{cBad}3}{\color{cBad}1}{\color{gray}4}$};
\node[font=\small] at (5.6,-1.18) {${\color{cBad}2}{\color{cBad}3}{\color{gray}4}{\color{cBad}1}$};
\node[font=\small] at (6.9,-1.18) {${\color{cBad}2}{\color{cBad}4}{\color{cBad}1}{\color{gray}3}$};
\node[font=\small] at (8.2,-1.18) {${\color{cBad}2}{\color{cBad}4}{\color{gray}3}{\color{cBad}1}$};
\fill[cA!12,rounded corners=2pt] (1.22,-1.87) rectangle (2.18,-1.49);
\node[font=\small] at (1.7,-1.68) {${\color{cA}\mathbf{3124}}$};
\node[font=\small] at (3,-1.68) {${\color{cBad}3}{\color{gray}1}{\color{cBad}4}{\color{cBad}2}$};
\fill[cA!12,rounded corners=2pt] (3.82,-1.87) rectangle (4.78,-1.49);
\node[font=\small] at (4.3,-1.68) {${\color{cA}\mathbf{3214}}$};
\node[font=\small] at (5.6,-1.68) {${\color{cBad}3}{\color{gray}2}{\color{cBad}4}{\color{cBad}1}$};
\node[font=\small] at (6.9,-1.68) {${\color{cBad}3}{\color{cBad}4}{\color{cBad}1}{\color{gray}2}$};
\node[font=\small] at (8.2,-1.68) {${\color{cBad}3}{\color{cBad}4}{\color{cBad}2}{\color{gray}1}$};
\fill[cA!12,rounded corners=2pt] (1.22,-2.37) rectangle (2.18,-1.99);
\node[font=\small] at (1.7,-2.18) {${\color{cA}\mathbf{4123}}$};
\fill[cA!12,rounded corners=2pt] (2.52,-2.37) rectangle (3.48,-1.99);
\node[font=\small] at (3,-2.18) {${\color{cA}\mathbf{4132}}$};
\fill[cA!12,rounded corners=2pt] (3.82,-2.37) rectangle (4.78,-1.99);
\node[font=\small] at (4.3,-2.18) {${\color{cA}\mathbf{4213}}$};
\node[font=\small] at (5.6,-2.18) {${\color{gray}4}{\color{cBad}2}{\color{cBad}3}{\color{cBad}1}$};
\fill[cA!12,rounded corners=2pt] (6.42,-2.37) rectangle (7.38,-1.99);
\node[font=\small] at (6.9,-2.18) {${\color{cA}\mathbf{4312}}$};
\fill[cA!12,rounded corners=2pt] (7.72,-2.37) rectangle (8.68,-1.99);
\node[font=\small] at (8.2,-2.18) {${\color{cA}\mathbf{4321}}$};
\draw[gray!40,line width=0.4pt] (0,-0.34) -- (10.4,-0.34);
\node[font=\small,anchor=west] at (0,0) {$n=3$};
\node[font=\small,anchor=west] at (0,-1.43) {$n=4$};
\node[font=\small,anchor=west,text=cA] at (8.85,0) {5 個避開};
\node[font=\small,anchor=west,text=cA] at (8.85,-1.43) {14 個避開};
\end{tikzpicture}}

\newcommand{\FigMelevenYoungSYT}{%
\begin{tikzpicture}
\fill[cA!14] (1,-0.36) -- (1.36,-0.36) -- (1.36,0) -- (1,0) -- cycle;
\draw[gray!60!black,line width=0.6pt] (1,-0.36) -- (1.36,-0.36) -- (1.36,0) -- (1,0) -- cycle;
\fill[cA!14] (1.36,-0.36) -- (1.72,-0.36) -- (1.72,0) -- (1.36,0) -- cycle;
\draw[gray!60!black,line width=0.6pt] (1.36,-0.36) -- (1.72,-0.36) -- (1.72,0) -- (1.36,0) -- cycle;
\fill[cA!14] (1.72,-0.36) -- (2.08,-0.36) -- (2.08,0) -- (1.72,0) -- cycle;
\draw[gray!60!black,line width=0.6pt] (1.72,-0.36) -- (2.08,-0.36) -- (2.08,0) -- (1.72,0) -- cycle;
\fill[cA!14] (2.08,-0.36) -- (2.44,-0.36) -- (2.44,0) -- (2.08,0) -- cycle;
\draw[gray!60!black,line width=0.6pt] (2.08,-0.36) -- (2.44,-0.36) -- (2.44,0) -- (2.08,0) -- cycle;
\fill[cA!14] (2.44,-0.36) -- (2.8,-0.36) -- (2.8,0) -- (2.44,0) -- cycle;
\draw[gray!60!black,line width=0.6pt] (2.44,-0.36) -- (2.8,-0.36) -- (2.8,0) -- (2.44,0) -- cycle;
\fill[cA!14] (1,-0.72) -- (1.36,-0.72) -- (1.36,-0.36) -- (1,-0.36) -- cycle;
\draw[gray!60!black,line width=0.6pt] (1,-0.72) -- (1.36,-0.72) -- (1.36,-0.36) -- (1,-0.36) -- cycle;
\fill[cA!14] (1.36,-0.72) -- (1.72,-0.72) -- (1.72,-0.36) -- (1.36,-0.36) -- cycle;
\draw[gray!60!black,line width=0.6pt] (1.36,-0.72) -- (1.72,-0.72) -- (1.72,-0.36) -- (1.36,-0.36) -- cycle;
\fill[cA!14] (1.72,-0.72) -- (2.08,-0.72) -- (2.08,-0.36) -- (1.72,-0.36) -- cycle;
\draw[gray!60!black,line width=0.6pt] (1.72,-0.72) -- (2.08,-0.72) -- (2.08,-0.36) -- (1.72,-0.36) -- cycle;
\fill[cA!14] (1,-1.08) -- (1.36,-1.08) -- (1.36,-0.72) -- (1,-0.72) -- cycle;
\draw[gray!60!black,line width=0.6pt] (1,-1.08) -- (1.36,-1.08) -- (1.36,-0.72) -- (1,-0.72) -- cycle;
\fill[cA!14] (1.36,-1.08) -- (1.72,-1.08) -- (1.72,-0.72) -- (1.36,-0.72) -- cycle;
\draw[gray!60!black,line width=0.6pt] (1.36,-1.08) -- (1.72,-1.08) -- (1.72,-0.72) -- (1.36,-0.72) -- cycle;
\fill[cA!14] (1.72,-1.08) -- (2.08,-1.08) -- (2.08,-0.72) -- (1.72,-0.72) -- cycle;
\draw[gray!60!black,line width=0.6pt] (1.72,-1.08) -- (2.08,-1.08) -- (2.08,-0.72) -- (1.72,-0.72) -- cycle;
\fill[cA!14] (1,-1.44) -- (1.36,-1.44) -- (1.36,-1.08) -- (1,-1.08) -- cycle;
\draw[gray!60!black,line width=0.6pt] (1,-1.44) -- (1.36,-1.44) -- (1.36,-1.08) -- (1,-1.08) -- cycle;
\fill[cA!14] (1.36,-1.44) -- (1.72,-1.44) -- (1.72,-1.08) -- (1.36,-1.08) -- cycle;
\draw[gray!60!black,line width=0.6pt] (1.36,-1.44) -- (1.72,-1.44) -- (1.72,-1.08) -- (1.36,-1.08) -- cycle;
\fill[cA!14] (1,-1.8) -- (1.36,-1.8) -- (1.36,-1.44) -- (1,-1.44) -- cycle;
\draw[gray!60!black,line width=0.6pt] (1,-1.8) -- (1.36,-1.8) -- (1.36,-1.44) -- (1,-1.44) -- cycle;
\node[font=\footnotesize,text=gray!40!black] at (3,-0.18) {5};
\node[font=\footnotesize,text=gray!40!black] at (2.28,-0.54) {3};
\node[font=\footnotesize,text=gray!40!black] at (2.28,-0.9) {3};
\node[font=\footnotesize,text=gray!40!black] at (1.92,-1.26) {2};
\node[font=\footnotesize,text=gray!40!black] at (1.56,-1.62) {1};
\node[font=\small] at (1.9,-2.15) {$\lambda=(5,3,3,2,1)\vdash 14$};
\fill[cA!14] (4.75,-0.3) -- (5.05,-0.3) -- (5.05,0) -- (4.75,0) -- cycle;
\draw[gray!60!black,line width=0.5pt] (4.75,-0.3) -- (5.05,-0.3) -- (5.05,0) -- (4.75,0) -- cycle;
\fill[cA!14] (5.05,-0.3) -- (5.35,-0.3) -- (5.35,0) -- (5.05,0) -- cycle;
\draw[gray!60!black,line width=0.5pt] (5.05,-0.3) -- (5.35,-0.3) -- (5.35,0) -- (5.05,0) -- cycle;
\fill[cA!14] (5.35,-0.3) -- (5.65,-0.3) -- (5.65,0) -- (5.35,0) -- cycle;
\draw[gray!60!black,line width=0.5pt] (5.35,-0.3) -- (5.65,-0.3) -- (5.65,0) -- (5.35,0) -- cycle;
\fill[cA!14] (5.65,-0.3) -- (5.95,-0.3) -- (5.95,0) -- (5.65,0) -- cycle;
\draw[gray!60!black,line width=0.5pt] (5.65,-0.3) -- (5.95,-0.3) -- (5.95,0) -- (5.65,0) -- cycle;
\fill[cA!14] (5.95,-0.3) -- (6.25,-0.3) -- (6.25,0) -- (5.95,0) -- cycle;
\draw[gray!60!black,line width=0.5pt] (5.95,-0.3) -- (6.25,-0.3) -- (6.25,0) -- (5.95,0) -- cycle;
\node[font=\small] at (5.5,-2.15) {$(5)$};
\fill[cA!14] (6.67,-0.3) -- (6.97,-0.3) -- (6.97,0) -- (6.67,0) -- cycle;
\draw[gray!60!black,line width=0.5pt] (6.67,-0.3) -- (6.97,-0.3) -- (6.97,0) -- (6.67,0) -- cycle;
\fill[cA!14] (6.97,-0.3) -- (7.27,-0.3) -- (7.27,0) -- (6.97,0) -- cycle;
\draw[gray!60!black,line width=0.5pt] (6.97,-0.3) -- (7.27,-0.3) -- (7.27,0) -- (6.97,0) -- cycle;
\fill[cA!14] (7.27,-0.3) -- (7.57,-0.3) -- (7.57,0) -- (7.27,0) -- cycle;
\draw[gray!60!black,line width=0.5pt] (7.27,-0.3) -- (7.57,-0.3) -- (7.57,0) -- (7.27,0) -- cycle;
\fill[cA!14] (7.57,-0.3) -- (7.87,-0.3) -- (7.87,0) -- (7.57,0) -- cycle;
\draw[gray!60!black,line width=0.5pt] (7.57,-0.3) -- (7.87,-0.3) -- (7.87,0) -- (7.57,0) -- cycle;
\fill[cA!14] (6.67,-0.6) -- (6.97,-0.6) -- (6.97,-0.3) -- (6.67,-0.3) -- cycle;
\draw[gray!60!black,line width=0.5pt] (6.67,-0.6) -- (6.97,-0.6) -- (6.97,-0.3) -- (6.67,-0.3) -- cycle;
\node[font=\small] at (7.27,-2.15) {$(4,1)$};
\fill[cA!14] (8.29,-0.3) -- (8.59,-0.3) -- (8.59,0) -- (8.29,0) -- cycle;
\draw[gray!60!black,line width=0.5pt] (8.29,-0.3) -- (8.59,-0.3) -- (8.59,0) -- (8.29,0) -- cycle;
\fill[cA!14] (8.59,-0.3) -- (8.89,-0.3) -- (8.89,0) -- (8.59,0) -- cycle;
\draw[gray!60!black,line width=0.5pt] (8.59,-0.3) -- (8.89,-0.3) -- (8.89,0) -- (8.59,0) -- cycle;
\fill[cA!14] (8.89,-0.3) -- (9.19,-0.3) -- (9.19,0) -- (8.89,0) -- cycle;
\draw[gray!60!black,line width=0.5pt] (8.89,-0.3) -- (9.19,-0.3) -- (9.19,0) -- (8.89,0) -- cycle;
\fill[cA!14] (8.29,-0.6) -- (8.59,-0.6) -- (8.59,-0.3) -- (8.29,-0.3) -- cycle;
\draw[gray!60!black,line width=0.5pt] (8.29,-0.6) -- (8.59,-0.6) -- (8.59,-0.3) -- (8.29,-0.3) -- cycle;
\fill[cA!14] (8.59,-0.6) -- (8.89,-0.6) -- (8.89,-0.3) -- (8.59,-0.3) -- cycle;
\draw[gray!60!black,line width=0.5pt] (8.59,-0.6) -- (8.89,-0.6) -- (8.89,-0.3) -- (8.59,-0.3) -- cycle;
\node[font=\small] at (8.74,-2.15) {$(3,2)$};
\fill[cA!14] (9.695,-0.3) -- (9.995,-0.3) -- (9.995,0) -- (9.695,0) -- cycle;
\draw[gray!60!black,line width=0.5pt] (9.695,-0.3) -- (9.995,-0.3) -- (9.995,0) -- (9.695,0) -- cycle;
\fill[cA!14] (9.995,-0.3) -- (10.295,-0.3) -- (10.295,0) -- (9.995,0) -- cycle;
\draw[gray!60!black,line width=0.5pt] (9.995,-0.3) -- (10.295,-0.3) -- (10.295,0) -- (9.995,0) -- cycle;
\fill[cA!14] (10.295,-0.3) -- (10.595,-0.3) -- (10.595,0) -- (10.295,0) -- cycle;
\draw[gray!60!black,line width=0.5pt] (10.295,-0.3) -- (10.595,-0.3) -- (10.595,0) -- (10.295,0) -- cycle;
\fill[cA!14] (9.695,-0.6) -- (9.995,-0.6) -- (9.995,-0.3) -- (9.695,-0.3) -- cycle;
\draw[gray!60!black,line width=0.5pt] (9.695,-0.6) -- (9.995,-0.6) -- (9.995,-0.3) -- (9.695,-0.3) -- cycle;
\fill[cA!14] (9.695,-0.9) -- (9.995,-0.9) -- (9.995,-0.6) -- (9.695,-0.6) -- cycle;
\draw[gray!60!black,line width=0.5pt] (9.695,-0.9) -- (9.995,-0.9) -- (9.995,-0.6) -- (9.695,-0.6) -- cycle;
\node[font=\small] at (10.145,-2.15) {$(3,1,1)$};
\fill[cA!14] (11.335,-0.3) -- (11.635,-0.3) -- (11.635,0) -- (11.335,0) -- cycle;
\draw[gray!60!black,line width=0.5pt] (11.335,-0.3) -- (11.635,-0.3) -- (11.635,0) -- (11.335,0) -- cycle;
\fill[cA!14] (11.635,-0.3) -- (11.935,-0.3) -- (11.935,0) -- (11.635,0) -- cycle;
\draw[gray!60!black,line width=0.5pt] (11.635,-0.3) -- (11.935,-0.3) -- (11.935,0) -- (11.635,0) -- cycle;
\fill[cA!14] (11.335,-0.6) -- (11.635,-0.6) -- (11.635,-0.3) -- (11.335,-0.3) -- cycle;
\draw[gray!60!black,line width=0.5pt] (11.335,-0.6) -- (11.635,-0.6) -- (11.635,-0.3) -- (11.335,-0.3) -- cycle;
\fill[cA!14] (11.635,-0.6) -- (11.935,-0.6) -- (11.935,-0.3) -- (11.635,-0.3) -- cycle;
\draw[gray!60!black,line width=0.5pt] (11.635,-0.6) -- (11.935,-0.6) -- (11.935,-0.3) -- (11.635,-0.3) -- cycle;
\fill[cA!14] (11.335,-0.9) -- (11.635,-0.9) -- (11.635,-0.6) -- (11.335,-0.6) -- cycle;
\draw[gray!60!black,line width=0.5pt] (11.335,-0.9) -- (11.635,-0.9) -- (11.635,-0.6) -- (11.335,-0.6) -- cycle;
\node[font=\small] at (11.635,-2.15) {$(2,2,1)$};
\fill[cA!14] (12.97,-0.3) -- (13.27,-0.3) -- (13.27,0) -- (12.97,0) -- cycle;
\draw[gray!60!black,line width=0.5pt] (12.97,-0.3) -- (13.27,-0.3) -- (13.27,0) -- (12.97,0) -- cycle;
\fill[cA!14] (13.27,-0.3) -- (13.57,-0.3) -- (13.57,0) -- (13.27,0) -- cycle;
\draw[gray!60!black,line width=0.5pt] (13.27,-0.3) -- (13.57,-0.3) -- (13.57,0) -- (13.27,0) -- cycle;
\fill[cA!14] (12.97,-0.6) -- (13.27,-0.6) -- (13.27,-0.3) -- (12.97,-0.3) -- cycle;
\draw[gray!60!black,line width=0.5pt] (12.97,-0.6) -- (13.27,-0.6) -- (13.27,-0.3) -- (12.97,-0.3) -- cycle;
\fill[cA!14] (12.97,-0.9) -- (13.27,-0.9) -- (13.27,-0.6) -- (12.97,-0.6) -- cycle;
\draw[gray!60!black,line width=0.5pt] (12.97,-0.9) -- (13.27,-0.9) -- (13.27,-0.6) -- (12.97,-0.6) -- cycle;
\fill[cA!14] (12.97,-1.2) -- (13.27,-1.2) -- (13.27,-0.9) -- (12.97,-0.9) -- cycle;
\draw[gray!60!black,line width=0.5pt] (12.97,-1.2) -- (13.27,-1.2) -- (13.27,-0.9) -- (12.97,-0.9) -- cycle;
\node[font=\small] at (13.27,-2.15) {$(2,1,1,1)$};
\fill[cA!14] (15.045,-0.3) -- (15.345,-0.3) -- (15.345,0) -- (15.045,0) -- cycle;
\draw[gray!60!black,line width=0.5pt] (15.045,-0.3) -- (15.345,-0.3) -- (15.345,0) -- (15.045,0) -- cycle;
\fill[cA!14] (15.045,-0.6) -- (15.345,-0.6) -- (15.345,-0.3) -- (15.045,-0.3) -- cycle;
\draw[gray!60!black,line width=0.5pt] (15.045,-0.6) -- (15.345,-0.6) -- (15.345,-0.3) -- (15.045,-0.3) -- cycle;
\fill[cA!14] (15.045,-0.9) -- (15.345,-0.9) -- (15.345,-0.6) -- (15.045,-0.6) -- cycle;
\draw[gray!60!black,line width=0.5pt] (15.045,-0.9) -- (15.345,-0.9) -- (15.345,-0.6) -- (15.045,-0.6) -- cycle;
\fill[cA!14] (15.045,-1.2) -- (15.345,-1.2) -- (15.345,-0.9) -- (15.045,-0.9) -- cycle;
\draw[gray!60!black,line width=0.5pt] (15.045,-1.2) -- (15.345,-1.2) -- (15.345,-0.9) -- (15.045,-0.9) -- cycle;
\fill[cA!14] (15.045,-1.5) -- (15.345,-1.5) -- (15.345,-1.2) -- (15.045,-1.2) -- cycle;
\draw[gray!60!black,line width=0.5pt] (15.045,-1.5) -- (15.345,-1.5) -- (15.345,-1.2) -- (15.045,-1.2) -- cycle;
\node[font=\small] at (15.195,-2.15) {$(1,1,1,1,1)$};
\fill[cA!10] (0,-3.2) -- (0.5,-3.2) -- (0.5,-2.7) -- (0,-2.7) -- cycle;
\draw[gray!60!black,line width=0.6pt] (0,-3.2) -- (0.5,-3.2) -- (0.5,-2.7) -- (0,-2.7) -- cycle;
\node[font=\small] at (0.25,-2.95) {$1$};
\fill[cA!10] (0.5,-3.2) -- (1,-3.2) -- (1,-2.7) -- (0.5,-2.7) -- cycle;
\draw[gray!60!black,line width=0.6pt] (0.5,-3.2) -- (1,-3.2) -- (1,-2.7) -- (0.5,-2.7) -- cycle;
\node[font=\small] at (0.75,-2.95) {$2$};
\fill[cA!10] (1,-3.2) -- (1.5,-3.2) -- (1.5,-2.7) -- (1,-2.7) -- cycle;
\draw[gray!60!black,line width=0.6pt] (1,-3.2) -- (1.5,-3.2) -- (1.5,-2.7) -- (1,-2.7) -- cycle;
\node[font=\small] at (1.25,-2.95) {$4$};
\fill[cA!10] (0,-3.7) -- (0.5,-3.7) -- (0.5,-3.2) -- (0,-3.2) -- cycle;
\draw[gray!60!black,line width=0.6pt] (0,-3.7) -- (0.5,-3.7) -- (0.5,-3.2) -- (0,-3.2) -- cycle;
\node[font=\small] at (0.25,-3.45) {$3$};
\fill[cA!10] (0.5,-3.7) -- (1,-3.7) -- (1,-3.2) -- (0.5,-3.2) -- cycle;
\draw[gray!60!black,line width=0.6pt] (0.5,-3.7) -- (1,-3.7) -- (1,-3.2) -- (0.5,-3.2) -- cycle;
\node[font=\small] at (0.75,-3.45) {$5$};
\node[font=\small,text=cA] at (0.75,-4.03) {是 SYT};
\fill[gray!6] (2.3,-3.2) -- (2.8,-3.2) -- (2.8,-2.7) -- (2.3,-2.7) -- cycle;
\draw[gray!60!black,line width=0.6pt] (2.3,-3.2) -- (2.8,-3.2) -- (2.8,-2.7) -- (2.3,-2.7) -- cycle;
\node[font=\small] at (2.55,-2.95) {$1$};
\fill[cBad!25] (2.8,-3.2) -- (3.3,-3.2) -- (3.3,-2.7) -- (2.8,-2.7) -- cycle;
\draw[gray!60!black,line width=0.6pt] (2.8,-3.2) -- (3.3,-3.2) -- (3.3,-2.7) -- (2.8,-2.7) -- cycle;
\node[font=\small] at (3.05,-2.95) {$4$};
\fill[gray!6] (3.3,-3.2) -- (3.8,-3.2) -- (3.8,-2.7) -- (3.3,-2.7) -- cycle;
\draw[gray!60!black,line width=0.6pt] (3.3,-3.2) -- (3.8,-3.2) -- (3.8,-2.7) -- (3.3,-2.7) -- cycle;
\node[font=\small] at (3.55,-2.95) {$5$};
\fill[gray!6] (2.3,-3.7) -- (2.8,-3.7) -- (2.8,-3.2) -- (2.3,-3.2) -- cycle;
\draw[gray!60!black,line width=0.6pt] (2.3,-3.7) -- (2.8,-3.7) -- (2.8,-3.2) -- (2.3,-3.2) -- cycle;
\node[font=\small] at (2.55,-3.45) {$2$};
\fill[cBad!25] (2.8,-3.7) -- (3.3,-3.7) -- (3.3,-3.2) -- (2.8,-3.2) -- cycle;
\draw[gray!60!black,line width=0.6pt] (2.8,-3.7) -- (3.3,-3.7) -- (3.3,-3.2) -- (2.8,-3.2) -- cycle;
\node[font=\small] at (3.05,-3.45) {$3$};
\node[font=\small,text=cBad] at (3.05,-4.03) {第 2 行 $4>3$};
\fill[cA!5] (6.235,-3.12) -- (6.655,-3.12) -- (6.655,-2.7) -- (6.235,-2.7) -- cycle;
\draw[gray!60!black,line width=0.6pt] (6.235,-3.12) -- (6.655,-3.12) -- (6.655,-2.7) -- (6.235,-2.7) -- cycle;
\node[font=\small] at (6.445,-2.91) {$1$};
\fill[cA!5] (6.655,-3.12) -- (7.075,-3.12) -- (7.075,-2.7) -- (6.655,-2.7) -- cycle;
\draw[gray!60!black,line width=0.6pt] (6.655,-3.12) -- (7.075,-3.12) -- (7.075,-2.7) -- (6.655,-2.7) -- cycle;
\node[font=\small] at (6.865,-2.91) {$2$};
\fill[cA!5] (7.075,-3.12) -- (7.495,-3.12) -- (7.495,-2.7) -- (7.075,-2.7) -- cycle;
\draw[gray!60!black,line width=0.6pt] (7.075,-3.12) -- (7.495,-3.12) -- (7.495,-2.7) -- (7.075,-2.7) -- cycle;
\node[font=\small] at (7.285,-2.91) {$3$};
\fill[cA!5] (6.235,-3.54) -- (6.655,-3.54) -- (6.655,-3.12) -- (6.235,-3.12) -- cycle;
\draw[gray!60!black,line width=0.6pt] (6.235,-3.54) -- (6.655,-3.54) -- (6.655,-3.12) -- (6.235,-3.12) -- cycle;
\node[font=\small] at (6.445,-3.33) {$4$};
\fill[cA!5] (6.655,-3.54) -- (7.075,-3.54) -- (7.075,-3.12) -- (6.655,-3.12) -- cycle;
\draw[gray!60!black,line width=0.6pt] (6.655,-3.54) -- (7.075,-3.54) -- (7.075,-3.12) -- (6.655,-3.12) -- cycle;
\node[font=\small] at (6.865,-3.33) {$5$};
\fill[cA!14] (7.995,-3.12) -- (8.415,-3.12) -- (8.415,-2.7) -- (7.995,-2.7) -- cycle;
\draw[gray!60!black,line width=0.6pt] (7.995,-3.12) -- (8.415,-3.12) -- (8.415,-2.7) -- (7.995,-2.7) -- cycle;
\node[font=\small] at (8.205,-2.91) {$1$};
\fill[cA!14] (8.415,-3.12) -- (8.835,-3.12) -- (8.835,-2.7) -- (8.415,-2.7) -- cycle;
\draw[gray!60!black,line width=0.6pt] (8.415,-3.12) -- (8.835,-3.12) -- (8.835,-2.7) -- (8.415,-2.7) -- cycle;
\node[font=\small] at (8.625,-2.91) {$2$};
\fill[cA!14] (8.835,-3.12) -- (9.255,-3.12) -- (9.255,-2.7) -- (8.835,-2.7) -- cycle;
\draw[gray!60!black,line width=0.6pt] (8.835,-3.12) -- (9.255,-3.12) -- (9.255,-2.7) -- (8.835,-2.7) -- cycle;
\node[font=\small] at (9.045,-2.91) {$4$};
\fill[cA!14] (7.995,-3.54) -- (8.415,-3.54) -- (8.415,-3.12) -- (7.995,-3.12) -- cycle;
\draw[gray!60!black,line width=0.6pt] (7.995,-3.54) -- (8.415,-3.54) -- (8.415,-3.12) -- (7.995,-3.12) -- cycle;
\node[font=\small] at (8.205,-3.33) {$3$};
\fill[cA!14] (8.415,-3.54) -- (8.835,-3.54) -- (8.835,-3.12) -- (8.415,-3.12) -- cycle;
\draw[gray!60!black,line width=0.6pt] (8.415,-3.54) -- (8.835,-3.54) -- (8.835,-3.12) -- (8.415,-3.12) -- cycle;
\node[font=\small] at (8.625,-3.33) {$5$};
\fill[cA!5] (9.755,-3.12) -- (10.175,-3.12) -- (10.175,-2.7) -- (9.755,-2.7) -- cycle;
\draw[gray!60!black,line width=0.6pt] (9.755,-3.12) -- (10.175,-3.12) -- (10.175,-2.7) -- (9.755,-2.7) -- cycle;
\node[font=\small] at (9.965,-2.91) {$1$};
\fill[cA!5] (10.175,-3.12) -- (10.595,-3.12) -- (10.595,-2.7) -- (10.175,-2.7) -- cycle;
\draw[gray!60!black,line width=0.6pt] (10.175,-3.12) -- (10.595,-3.12) -- (10.595,-2.7) -- (10.175,-2.7) -- cycle;
\node[font=\small] at (10.385,-2.91) {$2$};
\fill[cA!5] (10.595,-3.12) -- (11.015,-3.12) -- (11.015,-2.7) -- (10.595,-2.7) -- cycle;
\draw[gray!60!black,line width=0.6pt] (10.595,-3.12) -- (11.015,-3.12) -- (11.015,-2.7) -- (10.595,-2.7) -- cycle;
\node[font=\small] at (10.805,-2.91) {$5$};
\fill[cA!5] (9.755,-3.54) -- (10.175,-3.54) -- (10.175,-3.12) -- (9.755,-3.12) -- cycle;
\draw[gray!60!black,line width=0.6pt] (9.755,-3.54) -- (10.175,-3.54) -- (10.175,-3.12) -- (9.755,-3.12) -- cycle;
\node[font=\small] at (9.965,-3.33) {$3$};
\fill[cA!5] (10.175,-3.54) -- (10.595,-3.54) -- (10.595,-3.12) -- (10.175,-3.12) -- cycle;
\draw[gray!60!black,line width=0.6pt] (10.175,-3.54) -- (10.595,-3.54) -- (10.595,-3.12) -- (10.175,-3.12) -- cycle;
\node[font=\small] at (10.385,-3.33) {$4$};
\fill[cA!5] (11.515,-3.12) -- (11.935,-3.12) -- (11.935,-2.7) -- (11.515,-2.7) -- cycle;
\draw[gray!60!black,line width=0.6pt] (11.515,-3.12) -- (11.935,-3.12) -- (11.935,-2.7) -- (11.515,-2.7) -- cycle;
\node[font=\small] at (11.725,-2.91) {$1$};
\fill[cA!5] (11.935,-3.12) -- (12.355,-3.12) -- (12.355,-2.7) -- (11.935,-2.7) -- cycle;
\draw[gray!60!black,line width=0.6pt] (11.935,-3.12) -- (12.355,-3.12) -- (12.355,-2.7) -- (11.935,-2.7) -- cycle;
\node[font=\small] at (12.145,-2.91) {$3$};
\fill[cA!5] (12.355,-3.12) -- (12.775,-3.12) -- (12.775,-2.7) -- (12.355,-2.7) -- cycle;
\draw[gray!60!black,line width=0.6pt] (12.355,-3.12) -- (12.775,-3.12) -- (12.775,-2.7) -- (12.355,-2.7) -- cycle;
\node[font=\small] at (12.565,-2.91) {$4$};
\fill[cA!5] (11.515,-3.54) -- (11.935,-3.54) -- (11.935,-3.12) -- (11.515,-3.12) -- cycle;
\draw[gray!60!black,line width=0.6pt] (11.515,-3.54) -- (11.935,-3.54) -- (11.935,-3.12) -- (11.515,-3.12) -- cycle;
\node[font=\small] at (11.725,-3.33) {$2$};
\fill[cA!5] (11.935,-3.54) -- (12.355,-3.54) -- (12.355,-3.12) -- (11.935,-3.12) -- cycle;
\draw[gray!60!black,line width=0.6pt] (11.935,-3.54) -- (12.355,-3.54) -- (12.355,-3.12) -- (11.935,-3.12) -- cycle;
\node[font=\small] at (12.145,-3.33) {$5$};
\fill[cA!5] (13.275,-3.12) -- (13.695,-3.12) -- (13.695,-2.7) -- (13.275,-2.7) -- cycle;
\draw[gray!60!black,line width=0.6pt] (13.275,-3.12) -- (13.695,-3.12) -- (13.695,-2.7) -- (13.275,-2.7) -- cycle;
\node[font=\small] at (13.485,-2.91) {$1$};
\fill[cA!5] (13.695,-3.12) -- (14.115,-3.12) -- (14.115,-2.7) -- (13.695,-2.7) -- cycle;
\draw[gray!60!black,line width=0.6pt] (13.695,-3.12) -- (14.115,-3.12) -- (14.115,-2.7) -- (13.695,-2.7) -- cycle;
\node[font=\small] at (13.905,-2.91) {$3$};
\fill[cA!5] (14.115,-3.12) -- (14.535,-3.12) -- (14.535,-2.7) -- (14.115,-2.7) -- cycle;
\draw[gray!60!black,line width=0.6pt] (14.115,-3.12) -- (14.535,-3.12) -- (14.535,-2.7) -- (14.115,-2.7) -- cycle;
\node[font=\small] at (14.325,-2.91) {$5$};
\fill[cA!5] (13.275,-3.54) -- (13.695,-3.54) -- (13.695,-3.12) -- (13.275,-3.12) -- cycle;
\draw[gray!60!black,line width=0.6pt] (13.275,-3.54) -- (13.695,-3.54) -- (13.695,-3.12) -- (13.275,-3.12) -- cycle;
\node[font=\small] at (13.485,-3.33) {$2$};
\fill[cA!5] (13.695,-3.54) -- (14.115,-3.54) -- (14.115,-3.12) -- (13.695,-3.12) -- cycle;
\draw[gray!60!black,line width=0.6pt] (13.695,-3.54) -- (14.115,-3.54) -- (14.115,-3.12) -- (13.695,-3.12) -- cycle;
\node[font=\small] at (13.905,-3.33) {$4$};
\node[font=\small] at (10.385,-4.03) {形狀 $(3,2)$ 的全部 SYT：$f^{(3,2)}=5$};
\draw[gray!35,line width=0.4pt] (4.25,0.1) -- (4.25,-4.18);
\end{tikzpicture}}

\newcommand{\FigMelevenWarn}{%
\begin{tikzpicture}
\fill[cA!14] (0,0) -- (0.36,0) -- (0.36,0.36) -- (0,0.36) -- cycle;
\draw[gray!60!black,line width=0.6pt] (0,0) -- (0.36,0) -- (0.36,0.36) -- (0,0.36) -- cycle;
\fill[cA!14] (0.36,0) -- (0.72,0) -- (0.72,0.36) -- (0.36,0.36) -- cycle;
\draw[gray!60!black,line width=0.6pt] (0.36,0) -- (0.72,0) -- (0.72,0.36) -- (0.36,0.36) -- cycle;
\fill[cA!14] (0.72,0) -- (1.08,0) -- (1.08,0.36) -- (0.72,0.36) -- cycle;
\draw[gray!60!black,line width=0.6pt] (0.72,0) -- (1.08,0) -- (1.08,0.36) -- (0.72,0.36) -- cycle;
\fill[cA!14] (0,-0.36) -- (0.36,-0.36) -- (0.36,0) -- (0,0) -- cycle;
\draw[gray!60!black,line width=0.6pt] (0,-0.36) -- (0.36,-0.36) -- (0.36,0) -- (0,0) -- cycle;
\fill[cA!14] (0.36,-0.36) -- (0.72,-0.36) -- (0.72,0) -- (0.36,0) -- cycle;
\draw[gray!60!black,line width=0.6pt] (0.36,-0.36) -- (0.72,-0.36) -- (0.72,0) -- (0.36,0) -- cycle;
\node[font=\small] at (0.54,-1.05) {$5=3+2$};
\node[font=\Large,text=cBad] at (1.75,0) {$\neq$};
\node[font=\small] at (3.05,-1.05) {$\{1,2,5\},\{3,4\}$};
\draw[gray!65,line width=0.5pt] (3.05,0) circle (0.55);
\draw[tri] (3.05,0.55) -- (3.573,0.17) -- (2.527,0.17) -- cycle;
\draw[chord] (3.373,-0.445) -- (2.727,-0.445);
\fill (3.05,0.55) circle (0.04);
\node[font=\scriptsize] at (3.05,0.79) {1};
\fill (3.573,0.17) circle (0.04);
\node[font=\scriptsize] at (3.801,0.244) {2};
\fill (3.373,-0.445) circle (0.04);
\node[font=\scriptsize] at (3.514,-0.639) {3};
\fill (2.727,-0.445) circle (0.04);
\node[font=\scriptsize] at (2.586,-0.639) {4};
\fill (2.527,0.17) circle (0.04);
\node[font=\scriptsize] at (2.299,0.244) {5};
\end{tikzpicture}}

% 附錄 A 第 3 段：A.9 堆疊與佇列排序（M09）、A.10 模式迴避（M10）、A.11 整數分割與 Young 表（M11）
% 圖由 mit/make_figs_part3.py 產生；文中的數字都在該檔的 check_numbers() 裡驗算。

% ---------------------------------------------------------------------
\subsection{堆疊與佇列排序}\label{sec:m09}

把 $1,2,\ldots,n$ 的一個排列 $w=w_1w_2\cdots w_n$ 從左到右一個一個送進一台機器，
希望它依序輸出 $1,2,\ldots,n$；機器中間只有一個暫放的容器 \mitref{L3}。哪些排列排得出來？

\textbf{堆疊（後進先出）}只能從頂端放入、拿出。規則：頂端正好是下一個要輸出的數就\textbf{彈出}，
否則把下一個輸入\textbf{推入}；輸入用完而頂端不對就失敗。每一步其實沒有選擇（頂端對時若改成推入，
新推入的數比它大，會把它永遠壓住），所以照規則做會失敗，就表示怎麼做都排不出來。
例：$4213$ 依序推入 4、2、1，彈出 1、2，推入 3，彈出 3、4，輸出 $1,2,3,4$。

\textbf{佇列＋插隊線（先進先出）}從尾端排入、從頭端出列；另有一條插隊線，讓正好輪到的數直接輸出。
佇列排序有幾種版本，本講義與動畫 M09 的規則是：佇列頭正好是下一個要輸出的數就\textbf{出列}；
否則，下一個輸入正好是它就走\textbf{插隊線}；再否則把它\textbf{排入}佇列尾；輸入用完而佇列頭不對就失敗。
這也不會錯過任何排法（佇列頭對時先出列不吃虧；佇列頭不對時，輪到的數若去排隊，就會卡在它後面）。
例：$231$ 依序排入 2、3，1 走插隊線，出列 2、3。兩種規則都由程式對 $n\le 6$ 窮舉所有走法核對過。

\begin{figure}[H]
\centering
\FigMnineStuck
\caption{卡住的那一刻。(a) 堆疊排 $231$：輸出 1 之後要輸出 \textcolor{cB}{2}，它卻被 \textcolor{cBad}{3} 壓在下面。
(b) 佇列＋插隊線排 $321$：1 走插隊線之後要輸出 \textcolor{cB}{2}，它卻排在 \textcolor{cBad}{3} 後面。}
\label{fig:m09-stuck}
\end{figure}

說 $w$ 含有 $231$，是指 $w$ 裡有三個數依序是「中、大、小」（不必相鄰）；含有 $321$ 則是「大、中、小」
（精確的定義見第~\ref{sec:m10} 節）。結果如下 \mitref{L3}：

\begin{thmbox}[可排序的排列]
\begin{enumerate}[leftmargin=1.6em]
  \item 用堆疊排得出來 $\iff$ $w$ 避開 $231$（不含 $231$）。
  \item 用佇列＋插隊線排得出來 $\iff$ $w$ 避開 $321$。
\end{enumerate}
兩種排得出來的排列都恰有 $C_n$ 個。例如 $n=3$ 的 6 個排列中，堆疊只排不出 $231$，
佇列＋插隊線只排不出 $321$，各有 $5=C_3$ 個；$n=4$ 時各有 $14=C_4$ 個。
\end{thmbox}

\begin{notebox}[為什麼]
\textbf{堆疊。}若 $b,c,a$ 依序出現且 $a<b<c$：$c$ 到的時候 $a$ 還沒來，$b$、$c$ 只能先後推入，
$c$ 壓住了必須先輸出的 $b$（圖~\ref{fig:m09-stuck}(a)）。反過來，若照規則做卡住了：要輸出的 $m$ 在堆疊裡，
緊壓在它上面的是某個 $x>m$。$x$ 推入時 $m$ 在頂端卻沒有彈出，可見那時要輸出的是更小的 $a$；
$a$ 不在 $m$ 下面（否則永遠出不來），只能還在輸入裡、在 $x$ 之後，於是 $m,x,a$ 是「中、大、小」。

\textbf{佇列。}若 $c,b,a$ 依序出現且 $c>b>a$：$c$、$b$ 到的時候 $a$ 還沒輸出，兩者都只能排隊，
$c$ 擋在必須先輸出的 $b$ 前面（圖~\ref{fig:m09-stuck}(b)）。反過來，若 $w$ 避開 $321$，
佇列由頭到尾永遠遞增：若 $x$ 排入時佇列裡已有 $y>x$，此刻要輸出的 $a<x$ 不在佇列裡
（否則它是佇列裡最小的，就在佇列頭，應該先出列），只能還在輸入裡、在 $x$ 之後，於是 $y,x,a$ 是 $321$。
佇列一直遞增，輸入用完時要輸出的數就在佇列頭，不會卡住。
\end{notebox}

\begin{keybox}[推入 $=$ U，彈出 $=$ D]
堆疊排序時，把推入記成 U、彈出記成 D。每個數恰好推入一次、彈出一次，而且彈出之前一定已經推入，
所以得到一條 Dyck 路徑，例如 $4213\mapsto$ UUUDDUDD。
反過來，從 Dyck 路徑可以讀回排列：第 $k$ 個 D 彈出的是 $k$，和它配對的 U 推入的就是同一個數。
所以堆疊排得出來的排列與 Dyck 路徑一一對應，個數是 $C_n$。
同一個數的推入與彈出連起來，正好是第~\ref{sec:m08} 節由 Dyck 路徑得到的不交叉配對；
貫穿的例子 UUDUDD 對應的排列是 $312$。佇列那一邊的個數見下一節末。
\end{keybox}

% ---------------------------------------------------------------------
\subsection{模式迴避}\label{sec:m10}

\textbf{定義} \mitref{L3}。設 $w$ 是 $1,\ldots,n$ 的排列，$\pi$ 是 $1,\ldots,k$ 的排列（$k\le n$）。
若能在 $w$ 裡挑出 $k$ 個數（保持原來的先後，不必相鄰），它們的大小順序和 $\pi$ 一樣，
就說 $w$ \textbf{含有}模式 $\pi$；否則說 $w$ \textbf{避開} $\pi$（$\pi$-avoiding）。

\textbf{例。}$w=41532$ 的第 1、3、4 個數是 $4,5,3$，大小順序是「中、大、小」，和 $231$ 一樣，
所以 $41532$ 含有 $231$。再看 $52143$：從 5 個數挑 3 個共有 $\binom{5}{3}=10$ 種挑法，
逐一檢查，得到的模式只有 $321,312,213,132$，沒有 $231$，所以 $52143$ 避開 $231$（它也能用堆疊排序）。
把排列畫成點圖（第 $i$ 個數畫在高度 $w_i$），含有 $\pi$ 就是能挑出 $k$ 個點，形狀和 $\pi$ 的點圖一樣（圖~\ref{fig:m10-plots}）。

\begin{figure}[H]
\centering
\FigMtenPlots
\caption{排列的點圖。$41532$ 的 \textcolor{cTop}{4, 5, 3} 和模式 $231$ 的形狀相同；$52143$ 挑不出這個形狀。}
\label{fig:m10-plots}
\end{figure}

\textbf{全部列出來。}$n=3$ 的 6 個排列中只有 $231$ 自己含有 $231$，5 個避開；
$n=4$ 的 24 個排列中 14 個避開（圖~\ref{fig:m10-grid}）。又是 $C_3=5$、$C_4=14$。

\begin{figure}[H]
\centering
\FigMtenGrid
\caption{$n=3$ 與 $n=4$ 的全部排列（字典序）。\textcolor{cA}{藍底}：避開 $231$；
灰字：含有 $231$，\textcolor{cBad}{紅字}是找到的第一個。}
\label{fig:m10-grid}
\end{figure}

\begin{thmbox}[定理]
堆疊可排序 $\iff$ 避開 $231$（第~\ref{sec:m09} 節）。而且對\textbf{每一個}長度 3 的模式 $\pi$，
$1,\ldots,n$ 的排列中避開 $\pi$ 的都恰有 $C_n$ 個 \mitref{L3}。
程式把 $123,132,213,231,312,321$ 逐一數出，六個模式的數字完全相同：
\begin{center}
\begin{tabular}{c|cccccc}
$n$ & 1 & 2 & 3 & 4 & 5 & 6\\
\hline
避開 $\pi$ 的個數 & 1 & 2 & 5 & 14 & 42 & 132
\end{tabular}
\end{center}
\end{thmbox}

\textbf{為什麼避開 $231$ 的有 $C_n$ 個：根＋左＋右。}看最大的數 $n$ 的位置，寫成 $w=L\,n\,R$。
若 $L$ 裡有某個數大於 $R$ 裡的某個數，這兩個數和 $n$ 就是「中、大、小」。所以
\[
\text{$w$ 避開 $231$}\iff\text{$L$ 的數全部小於 $R$ 的數，而且 $L$、$R$ 各自避開 $231$}.
\]
（反方向：$L$ 全部小於 $R$ 時，一組「中、大、小」不可能跨在 $L$、$R$ 兩邊。）
這時 $L$ 是 $1,\ldots,k$ 的排列、$R$ 是 $k+1,\ldots,n-1$ 的排列。$n$ 是\textcolor{cTop}{根}、
$L$ 是\textcolor{cLeft}{左}、$R$ 是\textcolor{cRight}{右}，個數是 $\sum_{k=0}^{n-1}C_kC_{n-1-k}$，
正是遞迴式 \eqref{eq:segner}。例：$\textcolor{cLeft}{21}\,\textcolor{cTop}{5}\,\textcolor{cRight}{34}$ 避開 $231$；
$\textcolor{cLeft}{41}\,\textcolor{cTop}{5}\,\textcolor{cRight}{32}$ 左邊的 $4$ 大於右邊的 $3$，$4,5,3$ 就是前面找到的那個 $231$。

\textbf{六個模式只有兩類。}把排列左右反轉（倒過來讀）或上下翻轉（每個數 $x$ 換成 $n+1-x$），
含有的模式也跟著反轉或翻轉：$231$ 反轉成 $132$、翻轉成 $213$、兩個都做成 $312$，所以這四個模式的個數相同。
$123$ 與 $321$ 互為反轉，是另一類；它們也是 $C_n$ 個（$321$ 就是第~\ref{sec:m09} 節的佇列＋插隊線），
但需要另一種一一對應，這裡不證明。

% ---------------------------------------------------------------------
\subsection{整數分割與 Young 表}\label{sec:m11}

L3 最後換到另一種「分割」\mitref{L3}。$n$ 的一個\textbf{整數分割}是把 $n$ 寫成由大到小排好的正整數之和：
$\lambda=(\lambda_1,\ldots,\lambda_k)$，$\lambda_1\ge\lambda_2\ge\cdots\ge\lambda_k>0$，$\lambda_1+\cdots+\lambda_k=n$，
記作 $\lambda\vdash n$。$5$ 的整數分割有 7 個：
$5,\ 4+1,\ 3+2,\ 3+1+1,\ 2+2+1,\ 2+1+1+1,\ 1+1+1+1+1$。
如果順序不同算不同的，叫做\textbf{合成}（composition）：$5$ 的合成有 $2^4=16$ 個
（5 個 1 排成一列，中間 4 個空隙各自決定切或不切），其中由大到小排好的 7 個正是 5 的整數分割。
所以 $3+2$ 是整數分割 $(3,2)$；$2+3$ 不是由大到小，只是合成 $(2,3)$。

\begin{notebox}[注意：兩種「分割」不是同一個東西]
\begin{minipage}[c]{0.67\linewidth}
這裡的\textbf{整數分割}只看每一份有多大。Chu 論文的分割（第 2 節）是\textbf{集合分割}：
把圓周上\textbf{有編號的點}分成區塊，例如 $\{1,2,5\},\{3,4\}$（圖~\ref{fig:paper} 的第 5 個）。
兩者只有一點關係：把集合分割的區塊大小由大到小排好，就得到一個整數分割。
Chu 的 $\vn$、$\vv$ 指定的正是這個整數分割，例如第 3 節原文例子 $\vn=(1,1)$、$\vv=(2,3)$
指定區塊大小 $3$ 與 $2$，也就是 $(3,2)\vdash 5$；但 Chu 數的是不交叉的集合分割本身（這個例子有 5 個）。
L1--4 也沒有發展一般的不交叉集合分割 \mitref{L1--4}。
\end{minipage}\hfill
\begin{minipage}[c]{0.31\linewidth}
\centering
\FigMelevenWarn
\end{minipage}
\end{notebox}

\textbf{Young 圖} \mitref{L3}。把 $\lambda$ 畫成第 $i$ 列有 $\lambda_i$ 個格子、每列靠左、
由上往下不變長的圖（英式畫法）。圖~\ref{fig:m11-young} 上排左邊是 $\lambda=(5,3,3,2,1)\vdash 14$，
右邊是 5 的 7 個整數分割。

\textbf{標準 Young 表（SYT）} \mitref{L3}。把 $1,2,\ldots,n$ 各填一次到形狀 $\lambda$ 的格子裡，
使\textbf{每一列由左到右遞增、每一行由上到下遞增}。形狀 $\lambda$ 的 SYT 個數記作 $f^\lambda$。
例如圖~\ref{fig:m11-young} 下排最左邊的表（第一列 $1,2,4$、第二列 $3,5$）是 SYT，它就是貫穿例子的 SYT 拿掉 6；
旁邊的表每一列都遞增，但第 2 行上面是 4、下面是 3，不是 SYT。
形狀 $(3,2)$ 的 SYT 全部列出來恰有 5 個：$f^{(3,2)}=5$。

\begin{figure}[H]
\centering
\FigMelevenYoungSYT
\caption{上排：$\lambda=(5,3,3,2,1)$ 的 Young 圖，共 14 格（小字是每列的長度）；右邊是 5 的 7 個整數分割。
下排：一個 SYT 與一個反例（\textcolor{cBad}{紅格}所在的第 2 行不遞增）；右邊是形狀 $(3,2)$ 的全部 5 個 SYT，
第 2 個（深色）就是最左邊那一個。}
\label{fig:m11-young}
\end{figure}

L4 的問題就是：$f^\lambda$ 怎麼算？第~\ref{sec:m12} 節會看到 $f^{(n,n)}=C_n$（例如 $f^{(3,3)}=C_3=5$），
第~\ref{sec:m13} 節的鉤長公式則給出一般的 $f^\lambda$。
